% !TeX spellcheck = en_US
\documentclass[french]{yLectureNote}

\title{Mecanique Analytique}
\subtitle{Physique}
\author{Paulhenry Saux}
\date{\today}
\yLanguage{Français}
%lionels.calmels@cemes.fr
\professor{M.Brut}
\usepackage{graphicx}%----pour mettre des images
\usepackage[utf8]{inputenc}%---encodage
\usepackage{geometry}%---pour modifier les tailles et mettre a4paper
%\usepackage{awesomebox}%---pour les boites d'exercices, de pbq et de croquis ---d\'esactiv\'e pour les TP de PC
\usepackage{tikz}%---pour deiffner + d\'ependance de chemfig
% \usepackage{tabularx}%---pour dimensionner automatiquement les tableaux avec variable X
\usepackage{awesomebox}%---Pour les boites info, danger et autres
\usepackage{menukeys}%---Pour deiffner les touches de Calculatrice
\usepackage{fancyhdr}%---pour les en-t\^ete personnalis\'ees
\usepackage{blindtext}%---pour les liens
\usepackage{hyperref}%---pour les liens (\`a mettre en dernier)
\usepackage{caption}%---pour la francisation de la l\'egende table vers Tableau
\usepackage{pifont}
\usepackage{array}%---pour les tableaux
\usepackage{lipsum}
\usepackage{yFlatTable}
\usepackage{multicol}
\usepackage{tabularx}
\usepackage{booktabs}
\newcommand{\Lim}[1]{\lim\limits_{\substack{#1}}\:}
\renewcommand{\vec}{\overrightarrow}
\newcommand{\N}[0]{\mathbb{N}}
\newcommand{\dd}{\mathrm{d}}
\newcommand{\norm}[1]{||\vec{#1}||}
\newcommand{\fo}{\psi(\vec{r},t)}
\newcommand{\foe}{\psi(\vec{r},t)\*}
\newcommand{\HH}{\hat{H}}
\newcommand{\hb}{\hbar}
\newcommand{\lap}{\nabla^2}
\newcommand{\lapcc}{\frac{\partial^2 }{\partial x^2}+\frac{\partial^2 }{\partial y^2}+\frac{\partial^2 }{\partial z^2}}
\newcommand{\E}{\vec{E}(M)}
\newcommand{\B}{\vec{B}(M)}
\newcommand{\V}{V(M)}
\newcommand{\er}{\vec{e_{\rho}}}
\newcommand{\ep}{\vec{e_{\varphi}}}
\newcommand{\pip}{\pi^+}
\newcommand{\pim}{\pi^-}
\newcommand{\mv}{\mathcal{V}}
\DeclareMathOperator{\Div}{Div}
\newcommand{\rot}{\vec{\mathrm{rot}}}
\newcommand{\grad}{\vec{\mathrm{grad}}}
\newcommand{\qa}{q_{\alpha}}
\newcommand{\qdb}{\dot{q_{\beta}}}
\newcommand{\qb}{q_{\beta}}
\setcounter{chapter}{1}
\begin{document}
\chapter{Calcul variationnel}
\section{Variations et Actions}
Parfois on veut minimiser non pas une fonction mais une fonctionnelle : $$\int L(f'(x), f(x))\dd x$$

\checkInfo{Problème}{Avec une fonction à une variable, pour trouver un maxima, on calcule simplement la dérivée.

Pour une fonction à plusieurs variables, on calcule le gradient.

On s'intersse maintenant aux fonctionelles.
}
\begin{definition}[Fonctionelle]
    Une fonctionelle prend une fonction entrée et donne un nombre réel en sortie
\end{definition}
Par exemple $$\int_x{x_i}^{x_f}L(x,f(x), f'(x))\dd x$$ est une fonctionnelle.

Par exemple, l'énergie potentielle d'un cable entre 2 poteau dont on modélise la corde par $y=f(x)$.
%todo corde.svg

%todo c2-1

Pour minimiser,on va calculer 
$$\tilde{f}_{f_0(x)+\delta f(x)} \simeq \int_xi^{x_f} \dots \delta f(x)\dd x$$

On calcule 

\begin{flalign*}
    \delta \tilde{f} &= \tilde{f}_{f_0(x)+\delta f(x)}- \tilde{f}_{f_0(x)}\\
    &= \int_{x_i}^{x_f} L(x, f(x)+\delta f, f'(x)+\delta f'(x))-L(x, f(x), f'(x))\dd x\\
    &= \int_{x_i}^{x_f} \left(\frac{\partial L}{\partial f'}\delta f'+\frac{\partial L}{\partial f}\delta f+\dots\right) \dd x\\
    &= \int_{x_i}^{x_f} \left(\frac{\dd }{\dd x}\left(\frac{\partial L}{\partial f'}\delta f\right)-\frac{\dd }{\dd x}\left(\frac{\partial L}{\partial f'}\right)\delta f+\frac{\partial L}{\partial f}\delta f\right) \dd x\\
    &= \int_{x_i}^{x_f} \frac{\dd }{\dd x}\left(\frac{\partial L}{\partial f'}\delta f\right)\dd x + \int_{x_i}^{x_f} \left(\frac{\partial L}{\partial f}-\frac{\dd }{\dd x}\frac{\partial L}{\partial f'}\right)\delta f \dd x
\end{flalign*}
%TODO À corriger
On retrouve alors l'équation d'Euler-lagrange
\begin{definition}[Action]
    Pour q une trajectoire, on a 
    $$S = \int L(q(t),\dot{q(t)},t)\dd t$$
\end{definition}
\section{Multiplicateurs de Lagrange et Contraintes}
\subsection{Multiplicateur de Lagrange}
\checkInfo{Rappel : Multiplicateurs de Lagrange}{
Pour trouver les extrema d'une fonction $f(x,y)$ sous une contrainte $g(x,y) = 0$, on utilise la méthode du multiplicateur de Lagrange.
\begin{enumerate}
    \item On introduit une fonction auxiliaire $f_{\lambda}(x, y) = f(x, y) + \lambda g(x, y)$, où $\lambda$ est le multiplicateur de Lagrange.
    \item On cherche les points critiques de $f_{\lambda}$ en résolvant le système :
    \[ \begin{cases}
    \frac{\partial f_{\lambda}}{\partial x} = \frac{\partial f}{\partial x} + \lambda \frac{\partial g}{\partial x} = 0 \\
    \frac{\partial f_{\lambda}}{\partial y} = \frac{\partial f}{\partial y} + \lambda \frac{\partial g}{\partial y} = 0
    \end{cases} \]
    Cela donne les points critiques $(x^*(\lambda), y^*(\lambda))$ qui dépendent de $\lambda$.
    \item On injecte ces points dans l'équation de contrainte $g(x^*(\lambda), y^*(\lambda)) = 0$ pour trouver la ou les valeurs de $\lambda = \lambda^*$.
    \item Les points $(x^*(\lambda^*), y^*(\lambda^*))$ sont les extrema de $f$ sous la contrainte $g=0$.
\end{enumerate}
Cette méthode se généralise à plus de variables et plus de contraintes.
}
On a une fonction de plusieurs variable $F$ On veut déterminer ses extremas sous la contrainte $f(\vec{x})$. Cela se produit en $\vec{x_0}$, alors 
$\vec{\nabla}F // \vec{\nabla}f \Rightarrow \vec{\nabla} = \lambda \vec{\nabla} f$

Si on a 2 contraintes, le sous espace est de dimension $n-2$. Ainsi, $\vec{\nabla}f_1, \vec{\nabla}f_2$ sont $\perp$ à l'espace tangentiel et $\vec{\nabla}F\perp$ à l'espace tengent, donc va \^etre une combinaison linéaire des 2.

\begin{theorem}[Gradient et contrainte dans le cas général]
    $\vec{\nabla}F = \sum_{i=1}^k \lambda_i \vec{\nabla}f_i(\vec{x_0})$
\end{theorem}
\subsection{Application au calcul variationnel}
On veut mnimiser $$S = \int_{t_i}^{t_f}L(q,\dot{q}, t)$$ mais avec la contrainte  intégrale $$\int_{t_i}^{t_f}f(q,\dot{q}, t)\dd t=0$$

On peut écrire
$$\tilde{S} =\int_{t_i}^{t_f} (L-\sum\lambda f) \dd t$$
\subsection{Contrainte locale}
C'est une contrainte que doit satisfaire S à chaque instant. Les multiplicateurs de Lagrange dépendent du temps :
$$\tilde{S} = \int_{t_i}^{t_f}(L-\sum_{i=1}^k \lambda(t)f)\dd t$$

Mais il vaut mieux utiliser les contraintes locales dès le début avec des contraintes scléronomes/hétéronomes.

\subsubsection{Exemple avec une roue qui glisse}
%todo schema roule
\explanation{2}{Il s'agit de la contrainte que l'on oublie au début et que l'on injecte ensuite avec le multiplicateur.}
\begin{align*}
\dot{X} - R\dot{\theta} &= 0,\\
X - R\theta &= \text{cste} = 0,\explain{2}{right}{0}{0.5}{}\\
U &= \frac12 m\dot{X}^{2} + \frac12 I\dot{\theta}^{2},\\
\text{(disque homogène)} \qquad I &= mR^{2},\\
V &= -mgX\sin(\alpha),\\
L &= U - V
= \frac12 m\dot{X}^{2} + \frac12 I\dot{\theta}^{2} + mgX\sin(\alpha)
- \lambda\,(X - R\theta),\\
\delta S
&= \int_{t_i}^{t_f}\left(
m\dot{X}\,\delta\dot{X}
+ I\dot{\theta}\,\delta\dot{\theta}
+ mg\sin(\alpha)\,\delta X
- \lambda\,\delta X
+ \lambda R\,\delta\theta
\right)\dd t,\\
S
&= \int_{t_i}^{t_f}\left(
\frac12 m\dot{X}^{2}
+ \frac12 I\dot{\theta}^{2}
+ mgX\sin(\alpha)
- \lambda\,(X-R\theta)
\right)\dd t,\\
\frac{\partial S}{\partial X(t)} = 0
&\Rightarrow m\ddot{X} - mg\sin(\alpha) + \lambda = 0,\\
\frac{\partial S}{\partial \theta(t)} = 0
&\Rightarrow I\ddot{\theta} - \lambda R = 0,\\
\frac{\partial S}{\partial \lambda(t)} = 0
&\Rightarrow X = R\theta.
\end{align*}
On a donc un système de 3 équations que l'on va pouvoir utiliser pour obtenir l'équation finale.
Comme $X=R\theta$, on a aussi $\ddot{X}=R\ddot{\theta}$, donc
\[
\lambda = \frac{I}{R}\ddot{\theta}=\frac{I}{R^2}\ddot{X}.
\]
En remplaçant dans la première équation :
\[
m\ddot{X}-mg\sin(\alpha)+\frac{I}{R^2}\ddot{X}=0,
\]
d'où
\[
\left(m+\frac{I}{R^2}\right)\ddot{X}=mg\sin(\alpha).
\]
Ainsi,
\[
\ddot{X}=\frac{mg\sin(\alpha)}{m+\frac{I}{R^2}}.
\]

Pour un disque homogène, $I=mR^2$, donc
\[
\ddot{X}=\frac{g\sin(\alpha)}{2}.
\]
\chapter{Variable cycliques, formule de Beltrami et théorème de Noether}
\section{Variables cycliques}
\subsection{Définition}
Soit un système avec $N$ coordonnées généralisées $q_1,\dots,q_N$, et supposons que le lagrangien
$L(\dot{q}_\alpha,q_\alpha,t)$ ne dépend pas de $q_c$.

L'équation d'Euler-Lagrange pour $q_c$ s'écrit alors :
\[
\frac{\dd}{\dd t}\left(\frac{\partial L}{\partial \dot{q}_c}\right)
-
\frac{\partial L}{\partial q_c}
= 0.
\]

Comme $L$ ne dépend pas de $q_c$, on a :
\[
\frac{\partial L}{\partial q_c}=0.
\]

Donc :
\[
\frac{\dd}{\dd t}\left(\frac{\partial L}{\partial \dot{q}_c}\right)=0,
\]
et par conséquent :
\[
\frac{\partial L}{\partial \dot{q}_c}=\text{cste}.
\]
\subsection{Exemples}
\subsubsection{Exemple 1}
Exemple : $x_1, x_2, x_3$ sont les coordonnées cartésiennes et $V(\vec{r})$ est indépendant de $x_1$.

On a alors :
\[
L=\frac{1}{2}m\left(\dot{x}_1^{\,2}+\dot{x}_2^{\,2}+\dot{x}_3^{\,2}\right)-V(\vec{r}).
\]
Comme $L$ ne dépend pas de $x_1$, la coordonnée $x_1$ est cyclique. Donc :
\[
\frac{\partial L}{\partial x_1}=0
\qquad \Rightarrow \qquad
\frac{\dd}{\dd t}\left(\frac{\partial L}{\partial \dot{x}_1}\right)=0.
\]

Ainsi :
\[
\frac{\partial L}{\partial \dot{x}_1}=m\dot{x}_1= p_1\text{cste},
\]
\warningInfo{Potentiel}{Quand un potentiel ne dépend d'une coordonnée, la composante de l'impulsion dans cette m\^eme direction est conservée.}
\subsubsection{Exemple 2}
\begin{theorem}[Énergie cinétique en sphérique]
    $$\frac{1}{2}m\left(\dot{r}^{\,2}+r^{2}\dot{\theta}^{\,2}+r^{2}\sin^{2}\theta\,\dot{\varphi}^{\,2}\right)$$
\end{theorem}
En coordonnées sphériques $(r,\theta,\varphi)$, si le potentiel $V(\vec{r})$ est indépendant de $\varphi$, alors
\[
L=\frac{1}{2}m\left(\dot{r}^{\,2}+r^{2}\dot{\theta}^{\,2}+r^{2}\sin^{2}\theta\,\dot{\varphi}^{\,2}\right)-V(\vec{r}).
\]
Comme $L$ ne dépend pas de $\varphi$, cette coordonnée est cyclique. On a donc
\[
\frac{\partial L}{\partial \dot{\varphi}}=\text{cste}
\qquad \Rightarrow \qquad
mr^{2}\sin^{2}\theta\,\dot{\varphi}=\text{cste}.
\]
\warningInfo{Potentiel invariant autour d'un axe}{Quand un potentiel est invariant autour d'un axe, la composante du moment cinétique selon cet axe est conservée}
\subsection{Formule de Beltrami}
Si L est indépendant de t\marginCheck{Si il est indépendant explicitement du temps, bien sur les variables d'une trajectoire dépendent toujours du temps.}
\[
\text{Si } \frac{\partial L}{\partial t}=0,\ \text{alors}\ 
\frac{\dd}{\dd t}\!\left(\sum_{\alpha=1}^{N}\dot q_\alpha
\frac{\partial L}{\partial \dot q_\alpha}-L\right)=0.
\]

En effet,
\begin{align*}
\frac{\dd}{\dd t}\!\left(\sum_{\alpha}\dot q_\alpha\frac{\partial L}{\partial \dot q_\alpha}-L\right)
&=
\sum_{\alpha}\ddot q_\alpha\frac{\partial L}{\partial \dot q_\alpha}
+\sum_{\alpha}\dot q_\alpha\frac{\dd}{\dd t}\!\left(\frac{\partial L}{\partial \dot q_\alpha}\right)
-\frac{\dd L}{\dd t} \\
&=
\sum_{\alpha}\dot q_\alpha\left[\frac{\dd}{\dd t}\!\left(\frac{\partial L}{\partial \dot q_\alpha}\right)-\frac{\partial L}{\partial q_\alpha}\right]
-\frac{\partial L}{\partial t}.
\end{align*}

Avec les équations d'Euler--Lagrange,
\[
\frac{\dd}{\dd t}\!\left(\frac{\partial L}{\partial \dot q_\alpha}\right)-\frac{\partial L}{\partial q_\alpha}=0,
\]
on obtient donc
\[
\frac{\dd}{\dd t}\!\left(\sum_{\alpha=1}^{N}\dot q_\alpha\frac{\partial L}{\partial \dot q_\alpha}-L\right)=0
\]
\begin{theorem}[Formule de Beltrami]
    Si L ne dépend pas du temps, 
    $$\sum_{\alpha=1}^{N}\dot q_\alpha\frac{\partial L}{\partial \dot q_\alpha}-L=\text{cste}$$
\end{theorem}
Cette constante vaut l'énergie mécanique. En effet, 
Pour une particule de masse $m$ soumise à un potentiel $V(\vec r)$,
\[
L=\frac12 m\dot{\vec r}^{\,2}-V(\vec r).
\]
Alors, d'après la formule de Beltrami,
\[
\sum_i \dot r_i\,\frac{\partial L}{\partial \dot r_i}-L
=\frac12 m\dot{\vec r}^{\,2}+V(\vec r)
=\text{cste}=E_{\mathrm m}.
\]

\subsection{Retour sur la corde}
Si une corde de masse linéique $\mu$ est décrite par $y=f(x)$, alors
\[
\dd l=\sqrt{\dd x^2+\dd y^2}
=\sqrt{1+\left(\frac{\dd y}{\dd x}\right)^2}\,\dd x.
\]
Comme $h=y(x)$ et $\dd m=\mu\,\dd l$, l'énergie potentielle vaut
\[
V=\int g\,h\,\dd m
=\mu g\int_{x_i}^{x_f} y(x)\sqrt{1+\left(y'(x)\right)^2}\,\dd x.
\]

On doit minimiser la fonction, en prenant en compte la contrainte :
\[
S = \int_{0}^{D} \mu g y \sqrt{1 + y'^2} + \lambda \left(\sqrt{1 + y'^2} - \frac{l}{D}\right) \dd x
\]
où $L$ est le Lagrangien.

Avec $\delta S = 0$, on a :
\[
\frac{\partial L}{\partial x} = 0 \quad \Rightarrow \quad y' - \frac{\partial L}{\partial y'} = \text{Cste}
\]

Ainsi :
\[
(\mu g y + \lambda) y'^2 - (\mu g y + \lambda) \sqrt{1 + y'^2} = \text{Cste}
\]

Soit :
\[
\frac{1}{\sqrt{1 + y'^2}} (\mu g y + \lambda) = \text{Cste}
\]

En posant $\lambda = \mu g (y_0 + S)$, on obtient :
\[
f = \frac{y_0 + S}{c} \quad \Rightarrow \quad \frac{\partial}{\partial y'}\sqrt{1 + y'^2} = 0
\]

Puis :
\[
z = \frac{y + S}{c}
\]

\end{document}

